\section{Введение}
\subsection*{Цель} \addcontentsline{toc}{subsection}{Цель}
Определить радиус линзы исходя из размеров колец Ньютона.

\subsection*{Задачи} \addcontentsline{toc}{subsection}{Задачи}
\begin{enumerate}
	\item Измерить диаметры 10-15 колец Ньютона. По итогам измерений построить график зависимости квадрата диаметра кольца от его номера $d_k^2(k)$. С помощью графика определить радиус кривизны $R$ поверхности линзы;
	\item Качественно исследовать влияние углового размера источника на интерференционную картину. Угловой размер источника можно менять, передвигая неоновую лампу;
	\item Провести наблюдение колец Ньютона в белом свете. Определить количество наблюдаемых цветных колец.  
\end{enumerate}


\subsection*{Приборы и оборудование} \addcontentsline{toc}{subsection}{Приборы и оборудование}
\begin{enumerate}
	\item Установка (микроскоп, микрометр, линза);
	\item Источник квантовомеханического света -- неоновая лампа;
	\item Источник белого света -- лампа накаливания;
\end{enumerate}